% GPT-6 Astra Pro; /root; 22 November 2026 UTC.
% New full proof; independent review required before promotion.
\section[HD-F1-S27: recovery of the cover and four nonzero trace-kernel classes]{HD-F1-S27: the recovered Hilbert curve recovers the actual cyclic cover, and the norm sheaf fails weak point-equivariant semiregularity}
\label{sec:f1-norm-fourfold-recovered-cover}

Work on the common nonempty open specified in S23 and S26.
Retain its actual connected etale triple cover
$\pi:Z\to C$, its Prym $P\subset JZ$, its restriction
polarization $\beta$, and the actual maps
$T\xrightarrow{b}Y\hookrightarrow P$.
Put $E=\mathcal O_Y$ and let $\mathfrak D_P$ be the
sixteen-dimensional polarized Weil tangent space of S23.
The exact map $\eta$ of S24 now satisfies
\[
 \ker(\eta|_{\mathfrak D_P})
       =d\operatorname{Pr}\bigl(H^1(C,T_C)\bigr),
       \qquad \operatorname{rank}(\eta|_{\mathfrak D_P})=4.
 \tag{F1S27.1}
\]
In particular the actual characteristic map induces an
injection
\[
 (\ker\mu_C)^*\lhook\joinrel\longrightarrow
 \ker\bigl(\sigma_E^2:\operatorname{Ext}^2_P(E,E)
                                      \to HH_{-2}(P)\bigr).
 \tag{F1S27.2}
\]
Its image is four dimensional. For every specified coherent
finite point-action lift of $E$, this entire image is invariant.
Consequently $E$ fails weak $G$-semiregularity for every such
lift, including the trivial action. This is a counterexample
for this actual norm-fourfold candidate on this open family;
the other perfect-object alternatives in the main question
remain open.

\paragraph{1. The actual Jacobian embedding supplied by the universal divisor.}
Let $D=\operatorname{Spec}\mathbf C[\epsilon]/(\epsilon^2)$,
and choose $\kappa\in\ker(\eta|_{\mathfrak D_P})$.
The exact lifting criterion in S24--S25 supplies an embedded
deformation $Y_D\subset P_D$ and its actual smooth canonical
resolution $T_D\to Y_D$. Here $(P_D,\beta_D,\omega_D)$ is
the polarized Weil deformation representing $\kappa$, with
$\omega_D^2+\omega_D+1=0$.
S26 constructs a smooth proper Hilbert curve $Z_D$ and
its universal relative divisor
$\mathcal D_D\subset Z_D\times_D T_D$.
It reduces to the specified marked divisors $D_z$.

Choose sections $z_0$ of $Z_D/D$ and $t_0$ of $T_D/D$;
smoothness lifts any specified central points. Normalize
$\mathcal O(\mathcal D_D)$ along these two sections.
The normalized line gives a morphism
\[
 Z_D\longrightarrow\operatorname{Pic}^0(T_D/D),
 \qquad z\longmapsto M_z\otimes M_{z_0}^{-1}.
\]
The relative Jacobian universal property extends it to an
actual homomorphism
\[
 \varphi_D:J_D:=J(Z_D/D)\longrightarrow
                     \operatorname{Pic}^0(T_D/D).
 \tag{F1S27.3}
\]
Relative Picard and Albanese varieties are abelian schemes
here: all the varieties are smooth projective over the
characteristic-zero Artin base, with the chosen sections.
The pullback induced by $a_D:T_D\to P_D$ is
\[
 a_D^*:P_D^\vee\longrightarrow\operatorname{Pic}^0(T_D/D).
 \tag{F1S27.4}
\]
Its central fibre is the actual isomorphism proved in S19
and S21. Hence (F1S27.4) is an isomorphism over $D$.
One can check this directly: a proper morphism with that
central fibre is finite and its algebra map is an isomorphism
by Nakayama and flatness of the two abelian schemes.
Dualizing (F1S27.4) gives an isomorphism
$A_D:\operatorname{Alb}(T_D/D)\to P_D$.

Use the same universal divisor in the other direction to
fix the sign and the integral map explicitly. Its finite
flat degree-eight projection to $T_D$ gives
\[
 u_D:T_D\longrightarrow J_D,\qquad
 t\longmapsto\mathcal O_{Z_D}(\mathcal D_t-\mathcal D_{t_0}).
\]
The relative Albanese universal property factors this actual
map through a homomorphism
$U_D:\operatorname{Alb}(T_D/D)\to J_D$.
Define
\[
                 i_D=U_DA_D^{-1}:P_D\longrightarrow J_D.
 \tag{F1S27.5}
\]
Its central fibre is literally the original inclusion
$i:P\hookrightarrow JZ$: the central universal divisor is
the specified degree-eight divisor in $\operatorname{Sym}^8Z$,
so $u_0(t)=i(a(t)-a(t_0))$ and $U_0=iA_0$.
This also specifies the transpose relation with (F1S27.3).
Under the convention $\lambda_L(x)=t_x^*L\otimes L^{-1}$,
point-divisor ratios give
$\varphi_D=-U_D^\vee\lambda_D$, where
$\lambda_D:J_D\xrightarrow{\sim}J_D^\vee$ is the principal
Jacobian polarization. The sign follows by restricting the
normalized universal divisor to the two Abel curves, or
from the normalized Poincare addition isomorphism.
The construction of $i_D$ uses $u_D$ directly and therefore
has neither a sign nor an isogeny ambiguity.

The morphism $i_D$ is a closed immersion. It is proper, and
its underlying central map is a closed immersion, so it is
quasi-finite and therefore finite. The map from the ambient
algebra to its finite pushforward is surjective modulo
$\epsilon$; Nakayama gives surjectivity over $D$.
Thus the Hilbert divisor has recovered the original Prym
factor as an actual abelian subscheme of the recovered Jacobian.

\paragraph{2. Rigidity and the exact polarized extension of the cyclic action.}
We use the following elementary rigidity fact in the
construction. Two homomorphisms between fixed abelian
schemes over $D$ which agree centrally agree over $D$.
For their difference reduces to zero, so its image lies
in an affine neighbourhood of the identity section.
Every function on the proper connected abelian scheme
comes from $D$, since $H^0(\mathcal O)=\mathbf C[\epsilon]/(\epsilon^2)$.
The map therefore factors through $D$; as a homomorphism
it takes the identity to the identity and is zero.
This also shows that the Hom group has no integer torsion:
a torsion homomorphism has image in a finite etale group,
and its connected source forces that homomorphism to be zero.

The induced polarization homomorphism
$\lambda_{P,D}=i_D^\vee\lambda_Di_D$ reduces to $\beta$.
The specified $\beta_D$ has the same central reduction,
so rigidity gives $\lambda_{P,D}=\beta_D$.
In the rational Hom group put
\[
 r_D=\beta_D^{-1}i_D^\vee\lambda_D,
 \qquad r_Di_D=1,
 \qquad s_D^{\mathrm{rat}}
       =1+i_D(\omega_D-1)r_D.
 \tag{F1S27.6}
\]
The projector $e_D=i_Dr_D$ is self-adjoint for $\lambda_D$.
Its image and kernel are orthogonal, and (F1S27.6) acts
as $\omega_D$ on its image and as the identity on its
kernel. Consequently
\[
 (s_D^{\mathrm{rat}})^3=1,
 \qquad (s_D^{\mathrm{rat}})^\vee\lambda_D
                        s_D^{\mathrm{rat}}=\lambda_D.
 \tag{F1S27.7}
\]
These identities also follow by multiplying (F1S27.6),
using $r_Di_D=1$ and $\omega_D^\vee\beta_D\omega_D=\beta_D$.
On the central fibre the orthogonal complement of $P$
is $\pi^*JC$, where the original deck transformation
$\sigma$ acts as the identity. It acts as $\omega$ on
$P$. Hence $s_0^{\mathrm{rat}}=\sigma$ on the original
Jacobian, through those two actual rational projectors.

We now prove that (F1S27.6) is integral. The polarization
type in S14 is $(1,1,1,1,3,3,3,3)$, so the kernel of
$\beta_D$ is killed by $3$. Its dual isogeny construction
gives an actual homomorphism $\beta_D':P_D^\vee\to P_D$
with $\beta_D'\beta_D=[3]$ and
$\beta_D\beta_D'=[3]$. Thus
$f_D=3s_D^{\mathrm{rat}}$ is an actual endomorphism of $J_D$.
On $J_D[3]$ it takes values in $J_D[3]$, since
$[3]f_D=f_D[3]$. Its central restriction is
$[3]\sigma|_{J[3]}=0$. The source and target are finite
etale over the local Artin base, so a morphism between
them is determined by its central fibre. Hence
$f_D|_{J_D[3]}=0$.
The quotient property of $[3]:J_D\to J_D$ factors it
uniquely as
\[
                 f_D=s_D[3]=[3]s_D,
              \qquad s_D\in\operatorname{End}_D(J_D).
 \tag{F1S27.8}
\]
This actual endomorphism has rationalization
$s_D^{\mathrm{rat}}$ and central value $\sigma$.
The torsion-freeness just proved turns (F1S27.7) into
the same integral identities for $s_D$. In particular
it is a polarization-preserving automorphism of order
three, with inverse $s_D^2$, and
$s_Di_D=i_D\omega_D$.

\paragraph{3. Infinitesimal Torelli recovers the actual cover deformation.}
Let $\xi\in H^1(Z,T_Z)$ be the marked first-order
class of $Z_D$, and let $d\tau(\xi)$ be its polarized
Jacobian class. The extension $s_D$ of $\sigma$ implies
\[
                   \sigma\cdot d\tau(\xi)=d\tau(\xi).
 \tag{F1S27.9}
\]
For example, choose local trivializations of a first-order
deformation and represent it by a tangent-valued Cech
cocycle. An extending automorphism is given locally by
its central automorphism plus an infinitesimal vector field.
Compatibility on overlaps says that the transformed cocycle
minus the old one is the Cech boundary of those vector
fields. This gives exactly (F1S27.9) for $J_D$.

The curve $Z$ is nonhyperelliptic. Besides the S13 input,
this follows directly from the free order-three action:
on a hyperelliptic curve of genus at least two the unique
degree-two quotient induces an order-three automorphism
of $\mathbf P^1$. Its two fixed points have fibres of
cardinality at most two, so the order-three transformation
fixes points above them, contradicting freeness.
An action trivial on the quotient would lie in its
order-two deck group and cannot have order three.

The infinitesimal Torelli map for this nonhyperelliptic
curve is injective. Its exact dual is
\[
 \operatorname{Sym}^2H^0(Z,K_Z)\longrightarrow H^0(Z,K_Z^2),
                         \qquad u\otimes v\longmapsto uv,
 \tag{F1S27.10}
\]
which is surjective by Max Noether's theorem.
The tangent formula is Griffiths, \emph{Some remarks and
examples on continuous systems and moduli}, J.~Math.
Mech.~16 (1967), Lemma3 and its proof, pp.~792--793.
For the whole published proof of the multiplication
theorem, use Berkesch--Schreyer, \emph{Syzygies, finite
length modules, and random curves}, MSRI Publ.~67 (2015),
the Noether Theorem2.1 and its proof, pp.~33--34.
The former calculation identifies the derivative using
the actual Abel differential; the relative-Jacobian
identification is also worked out in Landesman's
2019 expository note, Lemma2.1 and Theorem4.3.
The published multiplication result is used as stated.

All maps in (F1S27.10) are functorial for $\sigma$.
Thus injectivity and (F1S27.9) give $\sigma\xi=\xi$.
Since the original cover is etale,
$T_Z=\pi^*T_C$. Projection formula and the actual
averaging idempotent $(1+\sigma+\sigma^2)/3$ identify
\[
 \pi^*:H^1(C,T_C)\xrightarrow{\sim}H^1(Z,T_Z)^{\langle\sigma\rangle}.
 \tag{F1S27.11}
\]
Choose its unique inverse image $\xi_C$ of $\xi$.
An actual Cech cocycle for $\xi_C$ glues the local
trivial first-order deformations by $1+\epsilon\theta_{ij}$,
giving a smooth proper marked curve $C_D$.
The equivalence of finite etale covers across a nilpotent
thickening extends the original $\pi$ to
\[
                         \pi_D':Z_D'\longrightarrow C_D.
 \tag{F1S27.12}
\]
Its cyclic automorphism and group law extend uniquely
by full faithfulness. This is EGA IV, Theorem18.1.2;
the complete affine-lifting and gluing proof is Stacks,
Theorem41.15.2, Tag039R, with the finite assertion in
Lemma58.8.3, Tag0BQB.
Pulling back the displayed tangent cocycle shows that
$Z_D'$ has the class $\pi^*\xi_C=\xi$.
The Cech classification of first-order deformations
therefore supplies a marked isomorphism $Z_D'\simeq Z_D$.
It is unique because $H^0(Z,T_Z)=0$.
Transport (F1S27.12) along this isomorphism. We have
obtained an actual finite etale cyclic cover
$\pi_D:Z_D\to C_D$ extending the original one.

The Jacobian automorphism induced by its deck
transformation agrees centrally with $\sigma$.
Rigidity identifies it with the already constructed
$s_D$, so these two constructions retain exactly the
same cyclic action on the recovered Jacobian.

\paragraph{4. Equality with the original Prym deformation.}
We identify the Prym of (F1S27.12) with the specified
$P_D$, including its embedding, action and polarization.
Let $q_D:J_D\to J_D/i_D(P_D)$ be the actual abelian
quotient. The homomorphism $q_D(1-s_D)$ is zero on
the central fibre and hence is zero by rigidity.
Thus the image of $1-s_D$ lies in $i_D(P_D)$.
Conversely
\[
                  (1-s_D)i_D=i_D(1-\omega_D),
 \tag{F1S27.13}
\]
and $1-\omega_D$ is an isogeny, since
$(1-\omega_D)(1-\omega_D^2)=[3]$.
Its fppf image is all of $P_D$. Therefore the actual
image of $1-s_D$ is exactly $i_D(P_D)$.
For the finite etale cyclic cover this image is the
identity component of the kernel of the norm: the norm
annihilates $1-s_D$, and its differential is the trace
onto the five-dimensional invariant part, leaving an
eight-dimensional connected kernel. These maps and
dimensions identify the same abelian subscheme over $D$.
Equivalently, after the finite etale kernel components
are separated, the two smooth closed subgroups have the
same central fibre, and Nakayama gives equality.
Their polarizations agree by
$i_D^\vee\lambda_Di_D=\beta_D$, and their actions agree
by (F1S27.13).
Consequently the original marked tangent $\kappa$ equals
the Prym period tangent $d\operatorname{Pr}(\xi_C)$.

S23 had already shown that every actual cover tangent
belongs to $\ker\eta$. The preceding construction proves
the reverse inclusion on $\mathfrak D_P$ and hence
the equality in (F1S27.1). The period derivative has
rank twelve on this open by S23, while
$\dim\mathfrak D_P=16$. Its quotient has dimension four,
and (F1S27.1) gives rank four for $\eta$.

\paragraph{5. Four actual nonzero obstruction classes, with the finite action retained.}
Retain the actual maps of S24:
\[
 H^1(Y,N)\xrightarrow{j}\operatorname{Ext}^1_P(I_Y,E)
       \xrightarrow{d_2}\operatorname{Ext}^2_P(E,E),
       \qquad \operatorname{ob}_E=d_2j\eta.
 \tag{F1S27.14}
\]
Both $j$ and $d_2$ are injective there. Thus (F1S27.1)
makes the induced map from the specified period quotient
$(\ker\mu_C)^*$ injective. The exact full
trace-characteristic identity and entire character
calculation of S23 give
\[
              \sigma_E^2\operatorname{ob}_E(\kappa)
                       =\kappa\mathbin{\lrcorner}\operatorname{ch}(E)=0
                     \quad(\kappa\in\mathfrak D_P).
 \tag{F1S27.15}
\]
Equations (F1S27.14)--(F1S27.15) prove (F1S27.2).

Now take any specified coherent point lift, with each
element acting by $F_g=t_x^*(-)\otimes L$ for
$(x,L)\in P\times\widehat P$.
Translations preserve invariant forms and vector fields,
so act identically on $H^1(P,T_P)$. A line in
$\operatorname{Pic}^0(P)$ has zero Atiyah class: its
flat connection splits its first-jet extension.
The tensor identity
\[
 \operatorname{At}(F_gE)=t_x^*\operatorname{At}(E)
                    +1_E\otimes\operatorname{At}(L)
                       =t_x^*\operatorname{At}(E)
 \tag{F1S27.16}
\]
and naturality for the actual chosen lift
$F_gE\simeq E$ show that every class
$\operatorname{ob}_E(\kappa)$ is invariant under its
actual conjugation action on self-Ext.
This does not require a new choice of linearization.
The four-dimensional subspace in (F1S27.2) consequently
lies in the invariant kernel as well.
The identity-sector induction/forgetful comparison of
R05, with the same actual point functors and lift,
identifies this with a nonzero kernel of the weak
$G$-semiregularity map. Thus the asserted failure holds
for every such lift. The stabilizer bound $|G|\mid6$
from S19 remains valid, and no existence of a nontrivial
stabilizer is needed.

\paragraph{Framework and exact scope.}
The Hilbert component and its universal Cartier divisor
give the relative Picard homomorphism (F1S27.3).
The same universal divisor in the other direction gives
the literal inclusion (F1S27.5), with the stated transpose
for its Poincare biextension. The finite etale quotient and
nilpotent-lifting equivalence in EGA IV18.1.2 recover
the original cyclic cover. These are the Hilbert/Picard
and biextension maps used in FGA and SGA7I, VII--VIII;
the characteristic action, perfect trace and their
GRR/Adams compatibility remain those of S23--S24
and SGA6V in the established project conventions.
The variational Hodge tangent continues to annihilate
the whole character; (F1S27.14) now proves that four
of these directions obstruct this actual perfect object.
Accordingly this specific norm sheaf supplies no weakly
point-equivariant semiregular seed for DEG\_L
Q-T7-REPLACEMENT. The result asserts no nonexistence
of other perfect complexes or of other detecting cycles.
